% GENERATED FILE. Edit canonical lessons, then run npm run generate:books.
% canonical-source-hash: fnv1a64:8506b4e59ab0fa1b
% canonical-lessons: GU-C343-teko, GU-C343-himayat, GU-C343-viruddh, GU-C343-sahmat, GU-C343-dhalvu

\chapter{Taking Sides}
\label{ch:gu-taking-sides}
\hlchaptermodality{343}

\begin{chapteropening}
\textbf{By the end of this chapter:} \emph{I can say who is for an idea and who is against it.}

Complete the last lesson of chapter 343: I can say who is for an idea and who is against it.
\end{chapteropening}

\section[\texorpdfstring{\gu{ટેકો} (ṭeko)}{ટેકો (ṭeko)}]{\gu{ટેકો} (ṭeko) — backing; support for a person or a plan}

\label{lesson:GU-C343-teko}



\begin{quote}
Before the new one: say the Gujarati for a refutation, then the Gujarati for to cut someone short.
\end{quote}

\subsection*{You'll want to know: \gu{ટેકો}}
\textbf{\gu{ટેકો}} — \emph{ṭeko} — \textquotedblleft{}backing; support for a person or a plan\textquotedblright{}.

\begin{quote}
\gu{આખા} \gu{ગામે} \gu{આ} \gu{યોજનાને} \gu{ટેકો} \gu{આપ્યો}.
\end{quote}

\begin{quote}
\emph{ākhā gāme ā yojnāne ṭeko āpyo.}
\end{quote}

\begin{quote}
The whole village backed this plan.
\end{quote}

A masculine noun.

\begin{cousinweb}
\textbf{\gu{ટેકો}} (\emph{ṭeko}) CDIAL lists the related Gujarati \emph{ṭek}, \textquotedblleft{}support\textquotedblright{}, under a reconstructed \emph{ṭakk}, \textquotedblleft{}remain, stop\textquotedblright{} (CDIAL 5420): what keeps a thing standing.
\end{cousinweb}

\begin{grammarlens}[title={What you've built}]
One more word for giving reasons and opinions.
\end{grammarlens}

\subsection*{Your turn}
\begin{itemize}
\raggedright
  \item \emph{Say it:} \emph{ṭeko}
  \item \emph{Say it:} \emph{ṭeko}, once more
  \item \emph{Say it:} \emph{ākhā gāme ā yojnāne ṭeko āpyo.} — The whole village backed this plan.
  \item \emph{Recall:} say the Gujarati for a refutation, then the Gujarati for to cut someone short, then say \emph{ṭeko} again
\end{itemize}

\begin{tcolorbox}[breakable,colback=teal!4,colframe=teal!35!black,title={Before you move on}]
What does \gu{ટેકો} mean? (\textquotedblleft{}Backing; support for a person or a plan\textquotedblright{}.) Say it once more.
\end{tcolorbox}
\section[\texorpdfstring{\gu{હિમાયત} (himāyat)}{હિમાયત (himāyat)}]{\gu{હિમાયત} (himāyat) — advocacy}

\label{lesson:GU-C343-himayat}



\begin{quote}
Before the new one: say the Gujarati for to cut someone short, then the Gujarati for backing.
\end{quote}

\subsection*{You'll want to know: \gu{હિમાયત}}
\textbf{\gu{હિમાયત}} — \emph{himāyat} — \textquotedblleft{}advocacy\textquotedblright{}.

\begin{quote}
\gu{મીરા} \gu{છોકરીઓના} \gu{અભ્યાસની} \gu{હિમાયત} \gu{કરે} \gu{છે}.
\end{quote}

\begin{quote}
\emph{Mīrā chokrīonā abhyāsnī himāyat kare che.}
\end{quote}

\begin{quote}
Mira speaks up for girls' education.
\end{quote}

A feminine noun.

\begin{cousinweb}
\textbf{\gu{હિમાયત}} (\emph{himāyat}) is usually traced to Arabic \emph{ḥimāyat}, \textquotedblleft{}protection\textquotedblright{}, through Persian; Gujarati uses it for speaking up for a cause.
\end{cousinweb}

\begin{grammarlens}[title={What you've built}]
One more word for giving reasons and opinions.
\end{grammarlens}

\subsection*{Your turn}
\begin{itemize}
\raggedright
  \item \emph{Say it:} \emph{himāyat}
  \item \emph{Say it:} \emph{himāyat}, once more
  \item \emph{Say it:} \emph{Mīrā chokrīonā abhyāsnī himāyat kare che.} — Mira speaks up for girls' education.
  \item \emph{Recall:} say the Gujarati for to cut someone short, then the Gujarati for backing, then say \emph{himāyat} again
\end{itemize}

\begin{tcolorbox}[breakable,colback=teal!4,colframe=teal!35!black,title={Before you move on}]
What does \gu{હિમાયત} mean? (\textquotedblleft{}Advocacy\textquotedblright{}.) Say it once more.
\end{tcolorbox}
\section[\texorpdfstring{\gu{વિરુદ્ધ} (viruddh)}{વિરુદ્ધ (viruddh)}]{\gu{વિરુદ્ધ} (viruddh) — against}

\label{lesson:GU-C343-viruddh}



\begin{quote}
Before the new one: say the Gujarati for backing, then the Gujarati for advocacy.
\end{quote}

\subsection*{You'll want to know: \gu{વિરુદ્ધ}}
\textbf{\gu{વિરુદ્ધ}} — \emph{viruddh} — \textquotedblleft{}against\textquotedblright{}.

\begin{quote}
\gu{ઘણા} \gu{માણસો} \gu{આ} \gu{યોજનાની} \gu{વિરુદ્ધ} \gu{છે}.
\end{quote}

\begin{quote}
\emph{ghaṇā māṇso ā yojnānī viruddh che.}
\end{quote}

\begin{quote}
Many people are against this plan.
\end{quote}

\begin{cousinweb}
\textbf{\gu{વિરુદ્ધ}} (\emph{viruddh}) is a tatsama, borrowed from Sanskrit \emph{viruddha}, \textquotedblleft{}opposed, hindered\textquotedblright{} (Monier-Williams).
\end{cousinweb}

\begin{grammarlens}[title={What you've built}]
One more word for giving reasons and opinions.
\end{grammarlens}

\subsection*{Your turn}
\begin{itemize}
\raggedright
  \item \emph{Say it:} \emph{viruddh}
  \item \emph{Say it:} \emph{viruddh}, once more
  \item \emph{Say it:} \emph{ghaṇā māṇso ā yojnānī viruddh che.} — Many people are against this plan.
  \item \emph{Recall:} say the Gujarati for backing, then the Gujarati for advocacy, then say \emph{viruddh} again
\end{itemize}

\begin{tcolorbox}[breakable,colback=teal!4,colframe=teal!35!black,title={Before you move on}]
What does \gu{વિરુદ્ધ} mean? (\textquotedblleft{}Against\textquotedblright{}.) Say it once more.
\end{tcolorbox}
\section[\texorpdfstring{\gu{સહમત} (sahmat)}{સહમત (sahmat)}]{\gu{સહમત} (sahmat) — in agreement}

\label{lesson:GU-C343-sahmat}



\begin{quote}
Before the new one: say the Gujarati for advocacy, then the Gujarati for against.
\end{quote}

\subsection*{You'll want to know: \gu{સહમત}}
\textbf{\gu{સહમત}} — \emph{sahmat} — \textquotedblleft{}in agreement\textquotedblright{}.

\begin{quote}
\gu{હું} \gu{તમારા} \gu{મત} \gu{સાથે} \gu{સહમત} \gu{છું}.
\end{quote}

\begin{quote}
\emph{hũ tamārā mat sāthe sahmat chũ.}
\end{quote}

\begin{quote}
I agree with your view.
\end{quote}

\begin{cousinweb}
\textbf{\gu{સહમત}} (\emph{sahmat}) is built from Sanskrit \emph{saha}, \textquotedblleft{}with\textquotedblright{}, and \emph{mata}, \textquotedblleft{}an opinion\textquotedblright{} (Monier-Williams): of the same opinion.
\end{cousinweb}

\begin{grammarlens}[title={What you've built}]
One more word for giving reasons and opinions.
\end{grammarlens}

\subsection*{Your turn}
\begin{itemize}
\raggedright
  \item \emph{Say it:} \emph{sahmat}
  \item \emph{Say it:} \emph{sahmat}, once more
  \item \emph{Say it:} \emph{hũ tamārā mat sāthe sahmat chũ.} — I agree with your view.
  \item \emph{Recall:} say the Gujarati for advocacy, then the Gujarati for against, then say \emph{sahmat} again
\end{itemize}

\begin{tcolorbox}[breakable,colback=teal!4,colframe=teal!35!black,title={Before you move on}]
What does \gu{સહમત} mean? (\textquotedblleft{}In agreement\textquotedblright{}.) Say it once more.
\end{tcolorbox}
\section[\texorpdfstring{\gu{ઢળવું} (ḍhaḷvũ)}{ઢળવું (ḍhaḷvũ)}]{\gu{ઢળવું} (ḍhaḷvũ) — to lean towards; to slip down}

\label{lesson:GU-C343-dhalvu}



\begin{quote}
Before the new one: say the Gujarati for against, then the Gujarati for in agreement.
\end{quote}

\subsection*{You'll want to know: \gu{ઢળવું}}
\textbf{\gu{ઢળવું}} — \emph{ḍhaḷvũ} — \textquotedblleft{}to lean towards; to slip down\textquotedblright{}.

\begin{quote}
\gu{આખી} \gu{ચર્ચા} \gu{પછી} \gu{મારો} \gu{મત} \gu{તમારી} \gu{બાજુ} \gu{ઢળે} \gu{છે}.
\end{quote}

\begin{quote}
\emph{ākhī carcā pachī māro mat tamārī bāju ḍhaḷe che.}
\end{quote}

\begin{quote}
After the whole discussion my view leans your way.
\end{quote}

Past (it …): \gu{તે} \gu{દીવાલ} \gu{એક} \gu{બાજુ} \gu{ઢળી} \gu{ગઈ}, \emph{te dīvāl ek bāju ḍhaḷī gaī}.

\begin{cousinweb}
\textbf{\gu{ઢળવું}} (\emph{ḍhaḷvũ}) is inherited from \emph{ḍhalati}, \textquotedblleft{}bends over, falls\textquotedblright{}, a form CDIAL reconstructs; CDIAL glosses the Gujarati verb \textquotedblleft{}to slip, to be poured out\textquotedblright{} (CDIAL 5581). Leaning towards an opinion is a later sense.
\end{cousinweb}

\begin{grammarlens}[title={What you've built}]
One more word for giving reasons and opinions.
\end{grammarlens}

\subsection*{Your turn}
\begin{itemize}
\raggedright
  \item \emph{Say it:} \emph{ḍhaḷvũ}
  \item \emph{Say it:} \emph{ḍhaḷvũ}, once more
  \item \emph{Say it:} \emph{ākhī carcā pachī māro mat tamārī bāju ḍhaḷe che.} — After the whole discussion my view leans your way.
  \item \emph{Recall:} say the Gujarati for against, then the Gujarati for in agreement, then say \emph{ḍhaḷvũ} again
\end{itemize}

\begin{tcolorbox}[breakable,colback=teal!4,colframe=teal!35!black,title={Before you move on}]
What does \gu{ઢળવું} mean? (\textquotedblleft{}To lean towards; to slip down\textquotedblright{}.) Say it once more.
\end{tcolorbox}
